\documentclass[10pt,a4paper]{amsart}
\usepackage[margin=19mm]{geometry}
\usepackage{amsmath,amssymb,mathtools}
\usepackage[scr=rsfs,cal=euler]{mathalfa}
\usepackage{xfrac}
\usepackage[unicode,hidelinks,bookmarks=false]{hyperref}
\newcommand{\ie}{i.e.\ }
\newcommand{\cf}{cf.\ }
\newcommand{\resp}{respectively\ }
\newcommand{\Subsection}{\subsection}
\providecommand{\nobreakdash}{\nobreak\hbox{-}}
\setlength{\emergencystretch}{2em}
\multlinegap0pt
\usepackage{amssymb}
\usepackage{xcolor}
\newtheorem{lemma}{Lemma}
\renewcommand{\ge}{\geqslant}
\renewcommand{\le}{\leqslant}
\title{Biregular bipartite labeled multigraphs and perfect matchings in bipartite tensor products}
\author{Ilya I. Bogdanov, Fedor Petrov, Anton Sadovnichiy, Fedor Ushakov}
\date{Source: arXiv:2603.18253v1 (CC BY 4.0)}
\begin{document}
\maketitle

\noindent\emph{Source and licence.} Biregular bipartite labeled multigraphs and perfect matchings in bipartite tensor products. Version arXiv:2603.18253v1 (CC BY 4.0), distributed by arXiv under the Creative Commons Attribution 4.0 International licence, http://creativecommons.org/licenses/by/4.0/, as recorded in the arXiv licence metadata for this exact version and shown on the abstract page https://arxiv.org/abs/2603.18253v1. Full licence notice: \texttt{licenses/arxiv-2603.18253-CC-BY-4.0.txt}.\par\smallskip

\noindent\textit{Editorial scope.} Counted unit: the article's lemma find3 (source0.tex lines 348--350). The article's complete original proof of it is delivered with it (source0.tex lines 352--368, one \texttt{proof} environment of 17 lines). No mathematics is written, repaired or re-derived: the statement and the proof are the article's own lines, in the article's own notation, and no retained line is substituted or re-flowed.  No further source-local context is needed: every cross-reference of every retained line resolves inside the two delivered spans.  Nothing was omitted from any retained span, so the package carries no omission record.  No retained line contains a citation, so the wrapper carries no bibliography and no bibliography entry.  No retained line contains a figure environment, an image-inclusion command or drawing code, so no image is delivered and the package declares no asset dependency.  Editorial formatting only, in addition: an AMS \texttt{amsart} preamble that restates the article's own declarations and macros used by the retained lines, namely the environments \texttt{lemma} on separate counters, together with the standard packages those lines need.  The retained environments print in this document as Lemma 1 in order of appearance, whereas the article prints the same items with the numbering of its own full document.  The complete native source of the article as distributed in the arXiv e-print for this exact version is bundled unchanged as \texttt{sources/196682/source0.tex}.
\begin{lemma}\label{find3}
For a bipartite multigraph as described above, let $u$, $v$ and $f$ be three vertices from its left part such that $u$ and $v$ are not complementary. Then there exists a vertex $w$ in the left part, $w \notin \{u,v,f\}$, such that $|N(\{u,v,w\})| \ge 3m+3$ and no two vertices among $\{u,v,w\}$ are complementary.
\end{lemma}

\begin{proof}
Let $l = |N(\{u,v\})|$. We consider several cases.

\emph{Case 1:} There exist complementary vertices among the four vertices of the left part different from $u$ and $v$ (i.e., among the set consisting of $f$ and the three other vertices). Choose such a complementary pair and let $w$ be the vertex of that pair that is not $f$. Then $w$ is not complementary to $u$ or $v$ because each vertex has at most one complementary vertex. Moreover,
\[
|N(\{u,v,w\})| = |N(\{u,v\})| + |N(w)| \ge (2m+2) + (m+1) = 3m+3.
\]

\emph{Case 2:} Both $u$ and $v$ have complementary vertices (say $u'$ and $v'$, respectively). Then we can choose $w$ to be any vertex distinct from $u,v,f$ that is not complementary to $u$ or $v$ (such a vertex exists because there are six vertices in total and at most two complements). Then
\[
|N(\{u,v,w\})| \ge |N(u)| + |N(v)| + |N(w)| \ge (m+1) + (m+1) + (m+1) = 3m+3.
\]

\emph{Case 3:} Only one of $u$ and $v$, say $u$, has a complementary vertex. Then we can take two different vertices $w_1$ and $w_2$ from the left part, different from $u,v,f$, that are not complementary to $u$. Note that $l = |N(\{u,v\})| \ge (m+1) + |N(v)|$. The total number of edges from $w_1$ and $w_2$ to $N(\{u,v\})$ is at most $2|N(v)| - (2m+1) < 2(2|N(v)| - 2m - 1) \le 2(2l - 4m - 3)$. Hence for at least one of $w_1$ and $w_2$, say $w_1$, the number of edges to $N(\{u,v\})$ is less than $2l - 4m - 3$, so the number of edges to the remaining vertices of the right part is greater than $6m+4 - 2l$. Consequently, $|N(\{u,v,w_1\})| \ge 3m+3$.

\emph{Case 4:} Neither $u$ nor $v$ has a complementary vertex. Observe that there are at most two vertices $s$ different from $u$ and $v$ such that $|N(\{u,v,s\})| < 3m+3$. Therefore we can choose $w$ among the remaining vertices (excluding $f$) satisfying the required condition.
\end{proof}

\end{document}
